\chapter{Differential mathematics}
The importance of differential mathematics (the term I used, generally it is just \textit{differential equation} or such) should be emphasized, or should only be amplified, after the chapter on mathematical modelling. It is at here that its potential and interpretation is unrecognized, much to the unfortunate position of the chapter, but nevertheless, the topic is required to stand on its own in its part (you can't mix everything up well), and partake the comprehensive review purpose of further development which requires its knowledge. As such, we might as well copulate it right here, why differential equation and its subsequent theory is a thing, and the essence of it.

Differential equation and the like forms the corpus of mathematics, based on the newfound method of \textit{change}, or \textit{continuous change}. The term differential gives you the sense of the derivative, $dx$, and its traversal reverse operation, the integral $\int f(x)\: dx$. The calculus method and its essence of change, alongside with the development of the concept of density and continuity, makes way for describing the physical world, by its changing state. Suppose a system is characterized by its state $V$, and a subsequent of related measures, called the \textit{parameter} $\{x_{i}\}$. The changes control method then, induces, for such system, a way of choice to construct the functional relation and evolution-style change of the system, such that, for measures varies over time, with alternation of $V$ as its subsequent observations. There, we can relate each parameter together, such to induce functions between them, and more structure abiding such law of change. There are many types of such measure altogether, such is for example, proxy of certain properties of the system that is passive or resultant of others, or the one in charge of inducing such change, coming from certain functional perspective that is \textit{reversible} or not. This is called the \textbf{discrete state change structure}, and while it is good, it does not capture much of the real world scale, or at least of what we consider smoothness. Differential structures changed that, by using relevant concepts in the theory of analysis (or its lesser but well-known variation as calculus) to describe scaling smoothness. Hence, something called a fix in place of the discrete state structure. That is one important thing of differential equation. 

While there indeed are many more questions and properties that make differential system inherently useful pretty much far more powerful than what we had ever since its introduction (which is why we have been using it in physics for a long, long time now), we might as well reserve such question, and reserve such inquiries, while finding out about the theory itself. Subsequent sections of this chapter then, will explore the innate nature, to a given degree of my own ability to dissect such, and to extrapolate what is there and where to advance, in finding a way to make sense, and form our own structure of differential. 

\section{The differential equation}

A differential equation is an equation, where the unknown is a function and both the function and its derivatives may appear in the equation. This is to relates also, not just the parameter set, but also their changes, their evolution, the inquired law and structure of functional consequence. 

Differential equations are essential for a mathematical description of physics, as in physics the main subject is to predict certain physical setting and the world at large, how to move and what will it do. Some typical example includes those of classical physics of which is related to the parameter of time. 
\begin{enumerate}
    \item \textbf{Newton's law}: As a high-level formulation, the equation $f=ma$, with $a=d^{2}x/dt^{2}$, can give us the following differential equation system:
    $$m \frac{d^{2}\mathbf{x}}{dt^{2}}(t)=\mathbf{f}\left( t,\mathbf{x}(t),\frac{d\mathbf{x}}{dt}(t) \right)$$
    with $\mathbf{x}(t)$ the position vector of the particle at time $t$. This is a second order Ordinary Differential Equation (ODE), which takes into account of the acceleration over a vectorized 3-space. 
    
    \item \textbf{Wave equations}: A wave perturbation $u$ propagating in time $t$ and in $\mathbb{R}^{3}$, labelled by $\textbf{x}=(x,y,z)$ through the media with wave speed $v>0$ is $$\frac{\partial^{2}u}{\partial t^{2}}(t,\mathbf{x})=v^{2}\left( \frac{\partial^{2}u}{\partial x^{2}} (t,\mathbf{x})+\frac{\partial^{2}u}{\partial y^{2}}(t,\mathbf{x})+\frac{\partial^{2}u}{\partial z^{2}}\right)$$ This is a second order in time and space PDE. 
\end{enumerate}

Categorization of differential equations are fairly simple, usually done using the notion in analysis. The equation of Newton's is called ordinary differential equation (ODE) - the unknown function depends on a single independent variable, $t$, similar to be called as \textit{univariate differential equation} in calculus term. The equation of wave is called partial differential equation (PDE), where the unknown function depends on two or more independent variables, $t,x,y$ and $z$, and their partial derivative appears in the equations, such then is to be called the \textit{multivariate differential equation}. Those dependent variables form the \textit{expressed state} of the structure that the equation describes, and is subjected to variations of the independent variables. 

The order of a differential equation is then the highest derivative order that appears in the equation. Thereby, Newton's equation could be called as second order differential equation, since there's the term $d^2x(t)/dx^2$. The same notation applies for partial DE. To see this in action, we will start with differential equation of first order in the next section. 
\section{First Order Equations}

Of first-order equation, where only first derivatives are involved, we would like to have some expository with exponential and eigenfunctions. 

\subsection{The Exponential function}
Often, the solution to differential equations involves the exponential function. Consider a function $f(x)=a^{x}$, where $a>0$ is constant. The derivative of this function is then given by 
\begin{equation*}
    \begin{split}
        \frac{df}{dx} & = \lim_{h\to 0} \frac{a^{x+h}-a^{x}}{h}\\
        & = a^{x} \lim_{h\to 0} \frac{a^{h}-1}{h}\\
        & = \lambda a^{x}\\
        & = \lambda f(x)
    \end{split}
\end{equation*}
of which $lambda=f'(0)=\mathrm{const}$. So the derivative of an exponential function is a multiple of the original function. 

\begin{definition}[Exponential]
    $exp(x)=e^{x}$ is the unique function $F$ satisfying $f'(x)=f(x)$ and $f(0)=1$
\end{definition}

We write the inverse function as $\ln{(x)}$ or $\log{(x)}$. The importance of the exponential function lies in it being an eigenfunction of the differential operator. 
\begin{definition}[Eigenfunction]
    An eigenfunction under the differential operator is a functin whose functional form is unchanged by the operator. Only its magnitude is changed, i.e. 
    \begin{equation}
        \frac{df}{dx} = \lambda f
    \end{equation}
\end{definition}

\subsubsection{The Exponential Decay}

Dynamical system and physical system can be described by the same mathematical structure. In this case, we have linear, constant coefficients, differential equations describe these two situations. 

The exponential decay equation is a famous example in nuclear physics. For $N$ particles with decay constant $k>0$, we have the equation $$N'=-kN$$
The negative sign is because the number of original atom before decaying into another states is finite and decreasing over time with the decay constant. Specifically, for initial condition $N(0)=N_{0}$, is $$N(t)=N_{0}e^{-kt}$$
The differential equation step can be made using integrating factor method, since it is separable and linear. Hence, $$(N'+kN)e^{kt}=0\implies (e^{kt}N)' =0\implies e^{kt}N=c,c\in \mathbb{R}$$
The initial condition $N_{0}=N(0)=c$, so the solution of the IVP problem is $$N(t)=N_{0}e^{-kt}$$ as followed. 
The half-life of a radioactive substance is also relevant. Typically, this is defined as: $$N(\tau)=\frac{N(0)}{2}$$ as for the time required for the initial time atom counts to reduce by half. 
A radioactive material constant $k$ and half-life $\tau$ are related by the equation $$h\tau=\ln{(2)}$$
A radioactive material $N$ hence can be expressed in terms of half-life: 
$$N(t)=N_{0}2^{-t/r}$$


\subsection{Linear Differential Equations}
To start our discovery into differential equation, we start with the basic one - first order equations, and linear one. Further specification also includes dealing with the basic of \textit{constant coefficients} in the process, though, since non-constant coefficient would be generalized later on with the preliminary contexts from constant case. 

A first order linear differential equation \index{first-order differential equation} is defined as followed:
\begin{definition}[First order ODE]
    A first order ODE on the unknown $y$ is:
    $$y'(t)=f(t,y(t))$$
    Where $f$ is given and $y'=\frac{dy}{dt}$. The equation is \textit{linear} iff the source function $f$ is linear on its second argument, $$y'=a(t)y+b(t)$$ The linear equation has constant coefficients iff both $a$ and $b$ are constants. Otherwise, the equation has variable coefficient.
\end{definition}

We denote by $y:D\subset \mathbb{R}\to \mathbb{R}$ a real valued function $y$ defined on a domain $D$. Such function is solution of the differential equation iff the equal is satisfied for all values of the independent variable $t$ on $D$. 

We now state a precise formula for the solutions of constant coefficient linear equations. The proof relies on the chain rule for derivatives:
\begin{theorem}[Constant coefficient, LDE Solution]
    The linear differential equation $$y'=ay+b$$ with $a\neq 0$ and constant $b$ has infinitely many solutions: $$y(t)=ce^{at}-\frac{b}{a}$$
\end{theorem}
The proof of the theorem is pretty simple. Consider the case $b=0$, so $$y'=ay,a\in \mathbb{R}$$
Then $$y'=ay\implies \frac{y'}{y}=a\implies \ln{(\lvert y \rvert )}'=a\implies \ln(\lvert y \rvert )=at+c_{0}$$
Where $c_{0}\in\mathbb{R}$ is an arbitrary integration constant. We use the FTC theorem on the last step:
$$\int \ln{(\lvert y \rvert )}' \, dt=\ln{(\lvert y \rvert )} $$
Compute the exponential on both sides:
$$y(t)=ce^{at}$$
This is the solution of the differential equation in the case that $b=0$. The case $b\neq 0$ can be converted into the case above:

\begin{equation*}
    \begin{split}
        y' & = ay+b\\
        & = a\left( y+\frac{b}{a} \right)\\
        \end{split}
\end{equation*}


This implies that $$\left( y+\frac{b}{a} \right)'=a\left( y+\frac{b}{a} \right)$$
since $(b/a)'=0$. Denoting $\tilde{y}=y+(b/a)$, the equation above is $\tilde{y}'=a \tilde{y}$. The solution to this equation is hence: $$y(t)=ce^{at}-\frac{b}{a}$$
This establishes the theorem. 

\subsubsection{Remark}

We solve the differential equation above $y'=ay$ by transforming it into a total derivative. This is clearly transparent in the calculation:
$$\ln{(\lvert y \rvert )}'=a\implies(\ln{(\lvert y \rvert )}-at)'=0$$
We set $\psi(t,y(t))'=0$ with $\psi(t,y(t))=\ln{(\lvert y(t) \rvert)}-at$. This is called the potential function. The total derivative is simple to integrate, since $\psi'=0\implies \psi=c_{0}$ for $c_{0}\in \mathbb{R}$. So the solution to such problem is: $$y(t)=\pm e^{c_{0}+at}$$ and by denoting $c=e^{c_{0}}$, we gain $$y(t)=ce^{at}$$
In the case $b\neq 0$, we have $$\displaystyle{\psi (t,y(t))=\ln{\left( \left\lvert  y(t) + \frac{b}{a}  \right\rvert \right)}-at}$$

\subsection{The Integrating Factor Method}
The argument of the theorem above (4.1.2) cannot be generalized to all linear equations with variable coefficient. Another way is proposed, and that is the integrating factor method (IFM). Using this method, we will prove the theorem. \index{integrating factor method}

\begin{proof}
    We write the equation as 
$$y'-ay=b$$
Then multiply the differential equation by a function $\mu$ called the integrating factor, such that $\mu y'-a\mu y=\mu b$. Then, choose a positive function $\mu$ such that $$-a\mu=\mu'$$
For any function $\mu$ solution on the above equation, the differential equation has the form: $$\mu y'+\mu'y=\mu b$$
But the LHS is the total derivative of a product of two functions: $$(\mu y)'=\mu b$$
This is the property we want in an integrating factor $\mu$. We want to find a function $\mu$ such that the LHS of the differential equation for $y$ can be written as a total derivative. The integrating function for this case is then $$\mu(t)=e^{-at}$$
The differential equation is the total derivative of a potential function $\psi(t,y)$, of which in this case is: $$\psi(t,y)=\left( y+\frac{b}{a} \right)e^{-at}$$ Notice that this potential function is the exponential of the potential function found in the first proof of this Theorem. The differential equation for y is a total derivative,
$$
\frac{d\psi}{dt} (t,y(t)) = 0
$$
so, it is simple to integrate 
\begin{equation*}
    \psi(t, y(t)) =  c \Rightarrow \left(y(t)+\frac{b}{a}\right) e{-at} = c \Rightarrow y(t) = ce^{at} - \frac{b}{a}
\end{equation*}
which establishes the theorem, as required. 
\end{proof}

\subsection{The Initial Value Problem}
An initial value problem means to find a solution to both a differential equation and an initial condition. 
\begin{definition}[IVP]
    The initial value problem (IVP) is to find all solutions $y$ to $$y'=ay+b$$ that satisfy the initial condition $$y(t_{0})=y_{0}$$ where $a,b,t_{0}$ and $y_{0}$ are given constant of such initial system. 
\end{definition}
The equation $y(t_{0})=y_{0}$ is called the initial condition of the problem, and the associated initial value problem of such differential function has a unique solution.
\begin{theorem}[IVP Solution, constant coefficient]
    Given the constant $a,b,t_{0},y_{0}\in\mathbb{R}$, with $a\neq 0$, the initial value problem $$y'-ay+b\quad y(t_{0})=y_{0}$$ has the unique solution $$y(t)=\left( y_{0}+\frac{b}{a} \right)e^{a(t-t_{0})}-\frac{b}{a}$$
\end{theorem}

\begin{proof}
    The general solution of the differential equation $y'=ay+b$ is given for any choice of the integration $c$: 
$$y(t)=ce^{at}-\frac{b}{a}$$
The initial condition determines the value of the constant $c$: $$y_{0}=y(t_{0})=ce^{at_{0}}-\frac{b}{a}\quad \iff \quad c=\left( y_{0}+\frac{b}{a} \right)e^{-at_{0}}$$
Substitute in, we have $$y(t)=\left( y_{0}+\frac{b}{a} \right)e^{a(t-t_{0})}-\frac{b}{a}$$
This establishes the theorem.
\end{proof}

\subsection{Linear Variable Coefficient Equation}\index{LVCE}

The idea of the previous section for solving the equation $y'=ay+b$ can be applied for variable cases. For example, consider $b=0$, $a(t)$ depends on $t$, we can solve this as:

\begin{equation}
    \frac{y'}{y}=a(t)\implies \ln{(\lvert y \rvert )}'=a(t)\implies \ln{(\lvert y(t) \rvert )}=A(t)+c_{0}
\end{equation}

Where $$A=\int a \, dt $$ is a primitive or antiderivative of $a$. 
The solution is hence: 
$$y=\pm e^{A(t)+c_{0}}=\pm e^{A(t)}e^{c_{0}}$$ so we get the solution $y(t)=ce^{A(t)}$ with $c=\pm e^{c_{0}}$. 

\begin{theorem}[Variable Coefficient, Solution]
    If the function $a,b$ are continuous, then $$y'=a(t)y+b(t)$$ Has infinitely many solutions given by $$y(t)=ce^{A(t)}+e^{A(t)}\int e^{-A(t)}b(t) \, dt $$ with $A(t)=\int a \, dt$, $c\in\mathbb{R}$. 
\end{theorem}
\begin{proof}
    We write the differential equation on one side: 
$$y'-ay=b$$
and then multiply the differential equation by a function of $\mu$, the integrating factor: 
$$\mu y'=a\mu y=\mu b$$
We need to find a function $\mu$ such that $-a\mu=\mu'$, hence $$\mu y'+\mu'y=\mu b\iff(\mu y)'=\mu b$$
We need to find such $\mu$ so that the LHS turns into the total derivative as above. Solving for $\mu'=-a\mu$ returns $$\ln{(\lvert \mu \rvert )}=-\int a \, dt+c_{0}=-A+c_{0}$$
This hence returns to $\mu=\pm e^{c_{0}}e^{-A}$. Since $c_{1}$ is a constant which will cancel out anyway, we have that $c_{0}=0$, hence $c_{1}=1$, the integrating factor is then $$u(t)=e^{-A(t)}$$
This function is an integrating factor. Substitute this into our equation: 
$$(e^{-A(t)}y)'=e^{-A}b$$
Integrating on both sides, we get: 
$$\psi(t,y)=c=(e^{-A}y)-\int e^{-A}b \, dt $$ being the potential function of the differential equation. The solution of the differential equation can be computed from the second equation, setting $\psi(t,y)=c$, then:
$$y(t)=ce^{A(t)}+e^{A(t)}\int e^{-A(t)}b(t) \, dt $$
This hence proves the theorem. 
\end{proof}

From here, we can also generalize the IVP problems for variable coefficient. 
\begin{definition}[IVP Solution, variable coefficient]
    Given continuous function $a,b$ with domain $(t_{1},t_{2})$, and constant $t_{0}\in(t_{1},t_{2})$, and $y_{0}\in\mathbb{R}$, the initial value problem with $y(t_{0})=y_{0}$ $$y'=a(t)y+b(t)$$ has the unique solution $y$ on the domain $(t_{1},t_{2})$ given by: $$y(t)=y_{0}c^{A(t)}+e^{A(t)}\int^{t}_{t_{0}} e^{-A(s)}b(s)  \, dt $$ With $A(t)=\int^{t}_{t_{0}}a(s) \, ds$ being a particular antiderivative of function $a$. 
\end{definition}
\begin{proof}
    Theorem 4.2.1 gives us the general solution for the above equation: 
$$y(t)=ce^{A(t)}+e^{A(t)}\int e^{-A(t)}b(t) \, dt,\quad\quad c\in \mathbb{R} $$
We will use the notation $$K(t)=\int e^{-A(t)}b(t) \, dt $$
and then introduce the initial condition in the equation, which fixes the constant $c$: 
$$y_{0}=y(t_{0})=ce^{A(t_{0})}+e^{A(t_{0})}K(t_{0})$$
So we get the constant $c$ by 
$$
c=y_{t_{0}}e^{-A(t_{0})}-K(t_{0})
$$
Using this expression in the general solution above then 
$$y(t)=y_{0}e^{A(t)-A(t_{0})}+e^{A(t)}(K(t)-K(t_{0}))$$
Let us introduce the particular primitives $\hat{A}(t)=A(t)-A(t_{0})$ and $\hat{K}(t)=K(t)-K(t_{0})$ which vanish at $t_{0}$, that is
$$\hat{A}(t)=\int ^{t}_{t_{0}}a(s) \, ds $$
and 
$$\hat{K}(t)=\int ^{t}_{t_{0}}e^{-A(s)}b(s) \, ds $$
Then the solution $y$ of the IVP has the form: 
\begin{equation*}
    \begin{split}
        y(t) & = y_{0}e^{\hat{A}(t)}+e^{A(t)}\int ^{t}_{t_{0}}e^{-A(s)}b(s) \, ds\\
        & = y_{0}e^{\hat{A}(t)}+e^{A(t)}\left( \frac{e^{A(t_{0})}}{e^{A(t_{0})}} \right)\int ^{t}_{t_{0}}e^{-A(s)}b(s) \, ds\\
        & = y_{0}e^{\hat{A}(t)}+e^{A(t)}( e^{A(t_{0})}e^{-A(t_{0})} )\int ^{t}_{t_{0}}e^{-A(s)}b(s) \, ds\\
        & = y_{0}e^{\hat{A}(t)}+e^{A(t)-A(t_{0})}\int ^{t}_{t_{0}}e^{-(A(s)-A(t_{0}))}b(s) \, ds\\
        & = y_{0}e^{\hat{A}(t)}+e^{\hat{A}(t)}\int ^{t}_{t_{0}}e^{-\hat{A}(s)}b(s) \, ds
        \end{split}
\end{equation*}
Once we rename the particular primitive $\hat{A}$ by $A$, we establish the theorem. 
\end{proof}
\subsection{The Bernoulli Equation}
Bernoulli differential equation \index{Bernoulli equation} is a first order nonlinear differential equation. We will study this type of differential equation during this introductory section, simply because of its application afterward. 

\begin{definition}
    The Bernoulli equation is $$y'=p(t)y+q(t)y^{n}$$
    where $p,q$ are given functions and $n\in \mathbb{R}$. 
\end{definition}

For $n\neq 0,1$, the equation is nonlinear as it has quadratic and up, terms. If $n=2$, we get the logistic equation: $$y'=ry\left( 1-\frac{y}{K} \right)$$
The Bernoulli equation is special in that it is a nonlinear equation that can be transformed into a linear equation. 

\begin{theorem}[Bernoulli Solutions]
    The function $y$ is a solution of the Bernoulli equation $$y'=p(t)y+q(t)y^{n}\quad n\neq 1$$
iff the function $v=1/y^{(n-1)}$ is solution of the linear differential equation $$v'=-(n-1)p(t)v-(n-1)q(t)$$
\end{theorem}
\begin{proof}
    Divide the Bernoulli equation by $y^{n}$, 
$$\frac{y'}{y^{n}}=\frac{p(t)}{y^{n-1}}+q(t)$$
Introduce the new unknown $v=y^{-(n-1)}$ and compute its derivative: 
$$v'=[y^{-(n-1)}]'=-(n-1)y^{-n}y'$$
which implies that $$-\frac{v'(t)}{n-1}=\frac{y'(t)}{y^{n}(t)}$$
If we substitute $v$ and this last equation into the Bernoulli equation we get: 
$$-\frac{v'}{n-1}=p(t)v+q(t)\implies v'=-(n-1)p(t)v-(n-1)q(t)$$
This establishes the Theorem.
\end{proof}
\section{Seperable Equations}
Separable equations can be solved just by integrating on both sides of the differential equation. We tried this idea to solve linear equation, it does not work, but it works for this kind. \index{separable equation(s)}
\begin{definition}
    A \textbf{separable} differential equation for the function $y$ is: $$h(y)y'=g(t)$$
Where $h$ and $g$ are given functions. 
\end{definition}
A separable differential equation has the following properties: 
\begin{enumerate}
    \item The left-hand side depends only on $y$. 
    \item The right-hand side depends only on $t$. 
    \item The left-hand side is of the form $(\text{something on}\: y)\times y'$. 
\end{enumerate}
We have the following theorem of the solution of separable equation (which I don't actually like... but whatever): 

\begin{theorem}[Separable Solutions]
    If $h$ and $g$ are continuous with $h\neq 0$, then $$h(y)y'=g(t)$$
Has infinitely many solution $y$ satisfying the algebraic equation 
$$H(y(t))=G(t)+c$$
Where $c\in\mathbb{R}$ is arbitrary, $H$ and $G$ are antiderivatives of $h$ and $g$.
\end{theorem}
An antiderivative of $h$ is $$H(x)=\int h(y) \, dy $$
while an antiderivative of $g$ is the function $G(x)=\int g(t) \, dt$. 
\begin{proof}
    Integrate with respect to $t$ on both side: 
$$\int h(y(t))y'(t) \, dt = \int g(t) \, dt + c  $$
Where $c$ is an arbitrary constant. Substitution method with $y'(t)dt=dy$, we have: 
$$\int h(y) \, dy = \int g(t) \, dt +c $$
To integrate on each side, it means to find a function $H$ and $G$ such that: $$H(y)=\int h(y) \, dy\quad G(x)=\int g(t) \, dt  $$
The equation then is established as: $$H(y)=G(t)+c$$
which implicitly defines a function $y$ depends on $t$. This establishes the theorem. 
\end{proof}
When the solution is given in terms of an algebraic equations, we say that the solution $y$ is given in implicit form. 
\begin{theorem}
    A function $y$ is a solution in \index{implicit form}\textbf{implicit form} of the equation $h(y)y'=g(t)$ iff the function $y$ is solution of the algebraic equation $$H(y(t))=G(t)+c$$
where $H$ and $G$ are any antiderivatives of $h$ and $g$. In the case that function $H$ is invertible, the solution $y$ above is given in explicit form iff is written as: $$y(t)=H^{-1}(G(t)+c)$$
\end{theorem}
\subsection{Euler Homogeneous Equations}
We are presented with the concept of homogeneous function. An example of an homogeneous function is the energy of a thermodynamic system, the energy $E$ of a fixed amount of gas is a function of the gas entropy, $S$ and the gas volume, $V$. Such energy is an homogeneous function of degree one, $$E(cS,cV)=cE(S,V)\quad \forall c\in\mathbb{R}$$

The most 'effective' homogeneous function to be studied at this stage of my learning, would be Euler's:

\begin{definition}[Euler]
    An Euler homogeneous differential equation has the form $$y'(t)=F\left( \frac{y(t)}{t} \right)$$
\end{definition}
Any function $F$ of $t,y$ that depends only on the quotient $y/t$ is scale invariant. This means that $F$ does not change when we do the transformation $y\to cy,t\to ct$, $$F\left( \frac{cy}{ct} \right)=F\left( \frac{y}{t} \right)$$

Scale invariant functions are a particular case of *homogeneous function of degree $n$*, which are function $f$ satisfying $$f(ct,cy)=c^{n}f(y,t)$$

More often than not, Euler homogeneous differential equation come from a differential equation $Ny'+M=0$, where both $N$ and $M$ are homogeneous function of the same degree. 
\begin{theorem}
    If the function $N$, $M$, of $t,y$ are homogeneous of the same degree, then the differential equation $$N(t,y)y'(t)+M(t,y)=0$$
is Euler homogeneous. 
\end{theorem}
To solve Euler homogeneous function, we will transform an Euler homogeneous equation into a separable equation, which we know how to solve.
\begin{theorem}
    The Euler homogeneous equation $$y'=F\left( \frac{y}{t} \right)$$
for the function $y$ determines a separable equation for $v=y/t$, given by $$\frac{v'}{(F(v)-v)}=\frac{1}{t}$$
\end{theorem}
The original homogeneous equation for the function $y$ is transformed into a separable equation for the unknown function $v=y/t$. One solves for $v$ in implicit or explicit form, and then transforms back to $y=tv$. 

\begin{proof}
    Introduce the function $v=y/t$ into the differential equation, $$y'=F(v)$$
We still need to replace $y'$ in terms of $v$. This is done as followed: $$y(t)=tv(t)\implies y'(t)=v'(t)+tv'(t)$$
Introducing these expressions into the differential equation for $y$, we get $$v+tv'=F(v)\implies \frac{v'}{(F(v)-v)}=\frac{1}{v}$$
The equation on the right is separable, hence establishes the theorem. 
\end{proof}
\section{Exact Differential Equation}
A differential equation is exact when it is a total derivative of a function, called potential function. Exact equations are simple to integrate - any potential function must be constant. The solution of the differential equation define level surfaces of a potential function. 

A semi-exact differential equation is an equation that is not exact but can be transformed into an exact equation after multiplication by a function, called an integrating factor.

\subsection{Exact Equation}
A differential equation is exact if certain parts of the differential equation have matching partial derivative. 

\begin{definition}
    An exact differential equation for $y$ is $$N(t,y)y'+M(t,y)=0$$
where the function $N,M$ satisfy: $$\partial_{t}N(t,y)=\partial_{y}M(t,y)$$
\end{definition}
The choice of notation here is shorten: $$\partial_{t}N(t,y)=\frac{\partial N(t,y)}{\partial t}$$
and $$\partial_{t}M(t,y)=\frac{\partial M(t,y)}{\partial t}$$
Here, the definition used $y$ both as a function, and an independent variable. 
\subsection{Exact Solution}
Exact differential equations can be rewritten as a total derivative of a function, called a potential function. 
\begin{theorem}
    If the differential equation $$N(t,y)y'+M(t,y)=0$$
is exact, then it can be written as: $$\frac{d\psi}{dt}(t,y(t))=0$$
where $\psi$ is the potential function and satisfies $N=\partial_{y}\psi,M=\partial_{t}\psi$. Therefore, the solution of the exact equation are given in implicit form as $$\psi(t,y(t))=c,c\in\mathbb{R}$$
\end{theorem}

The condition above, for which $\partial_{t}N=\partial_{y}M$  which guarantee the existence of a potential function, is called the Poincare theorem.

\begin{theorem}[Poincare]
    Continuously differentiable function $N,M$ on $t,y$ satisfy $$\partial_{t}N(t,y)=\partial_{y}M(t,y)$$
*iff* there is a twice continuously differentiable function $\psi$, depending on $t,y$ such that $$\partial_{y}\psi(t,y)=N(t,y), \quad \partial_{t}\psi(t,y)=M(t,y)$$
\end{theorem}
A differential equation defines the function $N$ and $M$. The exact condition is equivalent to the existence of $\psi$, related to $N$ and $M$. If we see it as the gradient, then $\nabla \psi=\langle M,N\rangle$. 
\subsection{Geometrical Interpretations}
An exact equation and its solution can be pictured on the graph of a potential function. This is called a **geometrical interpretation** of the exact equation. We saw that an exact equation $Ny'+M=0$ can be written as $$\frac{d\psi}{dt}=0$$
The solution is then $\psi(t,y(t))=c$, hence defining a level curves topology of the potential function. Given the level curve, $\mathbf{r}(t)=(t,y(t))\in \mathbb{R}^{ty}$ points to the level curve, while $\mathbf{r}'(t)=(1,y'(t))$ is tangent to the level curve. We then have that $$\mathbf{r}'\perp \nabla \psi \iff \mathbf{r}\cdot \nabla \psi=0\iff M+Ny'=0$$
The differential equation can be thought as the condition $\mathbf{r}'\perp \nabla \psi$. 

A potential function is also called a conserved quantity. This is yet another interpretation of the equation $d\psi/dt=0$, or its integral $\psi(t,y(t))=c$. If we call $c=\psi_{0}=\psi(0,y(0))$, the value of the potential function at the initial condition, then $\psi(t,y(t))=\psi_{0}$. 
\subsection{Semi-exact equations}
Sometimes a non-exact differential equation can be rewritten as an exact equation. One way this could happen is multiplying the differential equation by an appropriate function.
\begin{definition}[Semi-exact]
    A \textbf{semi-exact} differential equation is a non-exact equation that can be transformed into an exact equation after a multiplication by an integrating factor. 
\end{definition}

We generalize the notion of semi-exact differential equation to nonlinear differential equation:
\begin{theorem}
    If the equation $$N(t,y)y'+M(t,y)=0$$
is not \textbf{exact}, with $\partial_{t}N\neq \partial_{y}M$, then the function $N\neq 0$ and where the function $h$ defined as: 
$$h=\frac{\partial_{y}M-\partial_{t}N}{N}$$
depends only on $t$, not on $y$. Hence the equation below is exact: 
$$(e^{H}N)y'+ (e^{H}M)=0$$
where $H$ is an antiderivative of $h$: $$H(t)=\int h(t) \, dt $$
\end{theorem}

Note that The function $\mu(t)=e^{H(t)}$ is called an integrating factor. Also, any integrating factor $\mu$ is solution of the differential equation: $$\mu'(t)=h(t)\mu(t)$$ and multiplication by an integrating factor transforms a non-exact equation into an exact equation.
\subsection{Equation for inverse function}
Sometimes the equation for a function $y$ is neither exact nor semi-exact, but the equation for the inverse function $y^{-1}$ might be. We change the independent variable from $t$ to $x$. Therefore: 
$$N(x,y)y'=M(x,y)=0,\quad y=y(x)\quad y'=\frac{dy}{dx}$$
We denote by $x(y)$ the inverse of $y(x)$, that is $$x(y_{1})=x_{1}\iff y(x_{1})=y_{1}$$
Recall the identity relating derivative of a function and its inverse function: 
$$x'(y)=\frac{1}{y'(x)}$$
We need the following theorem: 
\begin{theorem}
    $Ny'+M=0$ is exact if and only if $Mx'+N=0$ is also exact. For semi-exact equations, there is a difference. 
\end{theorem}
\begin{proof}
    Write the differential equation of a function $y$ with values $y(x)$: 
$$N(x,y)y'+M(x,y)=0$$
and $\partial_{x}N=\partial_{y}(M)$.
If a solution $y$ is invertible, we denote $y^{-1}(y)=x(y)$, hence we have $$x'(y)=\frac{1}{y'(x(y))}$$
Divide the differential equation above by $y'$: 
$$N(x,y)+M(x,y)x'=0$$
Where now $y$ is the independent variable and the unknown function is $x$, with value $x(y)$, and the prime means $x'=dx/dy$. The condition for this last equation to be exact is $$\partial_{y}M=\partial_{x}N$$
which is exactly the same condition for the equation $Ny'+M=0$ to be exact. This establishes the theorem. 
\end{proof}
Sometimes, in the literature, the equation $Ny'+M=0$ and $N+Mx'=0$ is written together as followed:
$$N\: dy+M\:dx=0$$
The next result states a condition on the equation for the inverse function $x$ to be semi-exact. This condition is not equal to the condition on the equation for the function $y$ to be semi-exact.
\begin{theorem}
    If the equation $$Mx'+N=0$$
is not exact, with $\partial_{y}M\neq \partial_{x}N$, the function $M\neq 0$ and where the function $l$ defined as $$l=-\frac{\partial_{y}M-\partial_{x}N}{M}$$
depends only on $y$, not on $x$, then the equation below is exact $$(e^{L}M)x'+(e^{L}N)=0$$
where $L$ is an antiderivative of $l$, $$L(y)=\int l(y) \, dy $$
\end{theorem}
\section{Conclusion}

In this preliminaries section, we have formalized and discovered the first practical concept of differential equation, the introduction to the most basic, wide-spread, and effective type of differential equation - first order differential, and its variants under such consideration. There's more to differential equation, of which includes nonlinear equation. While we have explicit formula for LDE, there's none effective "everything" formula for nonlinear equations. This would be reserved for later sections to come. 